% Pista geo-23j2-q3 · Geometria
% Procedència: PAU Catalunya 2023 · Convocatòria de juny · Sèrie 5 · Qüestió 3
\documentclass[11pt,a4paper]{article}
\usepackage{pista}
\pistainfo{Catalunya 2023}{Convocatòria de juny}{Sèrie 5}

\begin{document}

\section*{\textcolor{blauPista}{Qüestió 3}}

\noindent Considereu les rectes a l'espai $r: x = -y = z + m$ i $s: \left.\begin{array}{r} x + y = 1 \\ x - z = 0 \end{array}\right\}$, en què $m$ és un paràmetre real.

\subsection*{\textit{a)} \normalfont Estudieu la posició relativa per als diferents valors del paràmetre $m$. \hspace*{\fill}\textnormal{[1,25~punts]}}

\pista{En primer lloc, escriu cada recta amb un paràmetre. En el cas de la recta $r$, iguala cada membre a $t$:
\[
x = t, \qquad -y = t, \qquad z + m = t
\]
i aïlla $x$, $y$, $z$. La recta $s$ ve donada per dues equacions: escriu simplement \mbox{$x = t$} (és un paràmetre lliure), i aïlla $y$ de la 1a equació i $z$ de la 2a equació.

\medskip
A continuació, compara els dos vectors directors. Si els vectors són paral·lels (o sigui, amb coordenades proporcionals), llavors les rectes són paral·leles o bé coincidents. Per decidir-ho, agafa un punt de $s$ i comprova si és de la recta $r$. Creus que hi ha algun valor de $m$ que ho faci possible?}

\subsection*{\textit{b)} \normalfont Calculeu $m$ perquè la distància entre les rectes $r$ i $s$ sigui de $\sqrt{2}$ unitats. \hspace*{\fill}\textnormal{[1,25~punts]}}

\pista{Si a l'apartat \textit{a)} t'han sortit rectes paral·leles, fes servir aquesta fórmula, amb un punt $P$ de la recta $r$, un punt $Q$ de la recta $s$ i el vector director $\vec{v}$:
\[
d(r, s) = \frac{|\vec{v} \times \vec{PQ}|}{|\vec{v}|}
\]
Com que el punt $P$ depèn de $m$, obtindràs una expressió amb la variable $m$.

\medskip
A continuació, iguala aquesta expressió a $\sqrt{2}$. Obtindràs una equació que sembla difícil però no ho és tant, perquè si eleves els dos costats al quadrat, llavors et quedarà una equació de segon grau amb la variable $m$.}

\end{document}
